\input{header}

\title{Writing the Intro, Background, \& Empirical Strategy}
\subtitle{ECON 692 -- Applied Economics Seminar}
\author{Andrew Hobbs}
\institute{University of San Francisco}
\date{April 9, 2026}

\begin{document}

\begin{frame}[plain]
  \titlepage
\end{frame}

% ══════════════════════════════════════════════════════════
%  Agenda & Check-in
% ══════════════════════════════════════════════════════════

\begin{frame}{Today's Agenda}
  \small
  \begin{tabular}{@{}r@{\hspace{8pt}}l@{\hspace{16pt}}l@{}}
    \toprule
    \textbf{Time} & \textbf{Duration} & \textbf{Activity} \\
    \midrule
    4:35 & 10 min & Check-in: where are you on the draft? \\
    4:45 & 25 min & Anatomy of intro, background, empirical strategy \\
    5:10 & 10 min & Break \\
    5:20 & 20 min & Good vs.\ bad examples \\
    5:40 & 25 min & GitHub Pages: project sites \& blogs \\
    6:05 & 85 min & Work sprint \\
    7:30 & 25 min & Peer feedback \\
    7:55 & 20 min & Wrap-up \& next steps \\
    \bottomrule
  \end{tabular}
\end{frame}

\begin{frame}{Check-in}
  Quick poll:

  \vspace{6pt}

  \begin{itemize}
    \item Who has a complete results section?
    \item Who has started writing prose (intro, background, anything)?
    \item Who is doing a blog or website instead of a traditional paper?
  \end{itemize}

  \vspace{6pt}

  \textbf{\textcolor{usfgreen}{Full Draft is due April 24}} --- two weeks from today.

  \vspace{4pt}

  That means a complete paper (or complete site): intro, background, empirical strategy, results, conclusion. Today we focus on the front half.
\end{frame}

% ══════════════════════════════════════════════════════════
\section{Anatomy of the Intro}
% ══════════════════════════════════════════════════════════

\begin{frame}{The Funnel Structure}
  Every good intro follows the same basic shape:

  \vspace{8pt}

  \begin{center}
  \begin{tikzpicture}[
    box/.style={draw=usfgreen, fill=usfgreen!10, rounded corners=3pt,
                minimum height=0.65cm, font=\small, align=center},
    arr/.style={->, thick, usfgreen!70}
  ]
    \node[box, minimum width=10cm] (hook) at (0,0) {Hook: why should anyone care?};
    \node[box, minimum width=7.5cm] (context) at (0,-1.1) {Context: what do we know so far?};
    \node[box, minimum width=5cm] (gap) at (0,-2.2) {Gap: what's missing?};
    \node[box, minimum width=3cm, fill=usfgold!20] (contribution) at (0,-3.3) {Your contribution};

    \draw[arr] (hook.south) -- (context.north);
    \draw[arr] (context.south) -- (gap.north);
    \draw[arr] (gap.south) -- (contribution.north);
  \end{tikzpicture}
  \end{center}

  \vspace{4pt}

  Your reader should know your research question by the end of paragraph one.
\end{frame}

\begin{frame}{What Goes in Each Paragraph?}
  \small
  \begin{columns}[T]
    \begin{column}{0.48\textwidth}
      \textbf{\textcolor{usfgreen}{Paragraph 1: The Hook}}
      \begin{itemize}\setlength\itemsep{0pt}
        \item A fact, a puzzle, or a policy question
        \item Why this matters \emph{now}
        \item Your research question (last sentence)
      \end{itemize}

      \vspace{4pt}

      \textbf{\textcolor{usfgreen}{Paragraphs 2--3: Context}}
      \begin{itemize}\setlength\itemsep{0pt}
        \item What the literature says
        \item What gap remains
      \end{itemize}
    \end{column}
    \begin{column}{0.48\textwidth}
      \textbf{\textcolor{usfgreen}{The Contribution Paragraph}}
      \begin{itemize}\setlength\itemsep{0pt}
        \item ``This paper estimates\ldots''
        \item What you do, what you find, why it matters
        \item One sentence per main result
      \end{itemize}

      \vspace{4pt}

      \textbf{\textcolor{usfgreen}{The Roadmap (last paragraph)}}
      \begin{itemize}\setlength\itemsep{0pt}
        \item ``Section 2 describes\ldots''
        \item Keep it short --- 3--4 sentences
      \end{itemize}
    \end{column}
  \end{columns}
\end{frame}

\begin{frame}{Common Intro Mistakes}
  \begin{itemize}
    \item \textbf{Too broad.} ``Since the beginning of time, economies have\ldots'' No. Start with your specific topic.

    \vspace{4pt}

    \item \textbf{Burying the contribution.} If your reader has to get to page 3 to find out what you did, you've lost them.

    \vspace{4pt}

    \item \textbf{No stakes.} Why should I care? Connect your question to something concrete --- a policy, a cost, a puzzle.

    \vspace{4pt}

    \item \textbf{Mixing intro and lit review.} The intro should cite a few key papers. The deep survey goes in the background section.
  \end{itemize}
\end{frame}

% ══════════════════════════════════════════════════════════
\section{Anatomy of the Background}
% ══════════════════════════════════════════════════════════

\begin{frame}{What the Background Section Does}
  The background section has one job: \textbf{position your paper}.

  \vspace{8pt}

  It's not a summary of everything you've read. It's an argument for why your paper needs to exist.

  \vspace{8pt}

  Organize around three buckets:

  \vspace{4pt}

  \begin{enumerate}
    \item \textbf{The phenomenon.} What do we know about the thing you're studying?
    \item \textbf{The method.} What approaches have others used? What are the limitations?
    \item \textbf{The gap.} What hasn't been done --- and why does that matter?
  \end{enumerate}

  \vspace{8pt}

  Your lit review should end with a sentence that makes your paper feel \emph{inevitable}.
\end{frame}

\begin{frame}{Common Background Mistakes}
  \begin{itemize}
    \item \textbf{The Wikipedia survey.} ``Smith (2010) found X. Jones (2012) found Y. Lee (2015) found Z.'' This isn't a literature review; it's a list.

    \vspace{4pt}

    \item \textbf{No connection to your contribution.} Every paper you cite should be there for a reason. If you can't explain why it's relevant to \emph{your} paper, cut it.

    \vspace{4pt}

    \item \textbf{Too long.} Two to three pages is plenty for a capstone. You're not writing a dissertation.

    \vspace{4pt}

    \item \textbf{No synthesis.} Group papers by theme, not by author. What do they collectively tell us? Where do they collectively fall short?
  \end{itemize}
\end{frame}

% ══════════════════════════════════════════════════════════
\section{Anatomy of the Empirical Strategy}
% ══════════════════════════════════════════════════════════

\begin{frame}{The Empirical Strategy Recipe}
  This section should convince a skeptic that your approach is sound.

  \vspace{8pt}

  Four ingredients, in order:

  \vspace{4pt}

  \begin{enumerate}
    \item \textbf{Data.} Unit of observation, time period, sample size, key variables. Where does it come from? What's excluded and why?
    \item \textbf{Identification.} What variation are you exploiting? Why is it plausibly exogenous? What's the identifying assumption?
    \item \textbf{Estimation.} The model. Write down the equation. Explain every term.
    \item \textbf{Inference.} How are standard errors clustered? Why that level?
  \end{enumerate}

  \vspace{6pt}

  If you don't have an equation in this section, something is missing.
\end{frame}

\begin{frame}{The Identification Argument}
  The identification argument is the heart of your empirical strategy. It answers: \textbf{why should I believe this is causal?}

  \vspace{6pt}

  \begin{columns}[T]
    \begin{column}{0.48\textwidth}
      \textbf{\textcolor{usfgreen}{For DiD:}}
      \begin{itemize}
        \item What's the treatment? When did it happen?
        \item Who's treated vs.\ control?
        \item Why is parallel trends plausible?
        \item What threatens it?
      \end{itemize}
    \end{column}
    \begin{column}{0.48\textwidth}
      \textbf{\textcolor{usfgreen}{For IV:}}
      \begin{itemize}
        \item What's the instrument?
        \item Why is it relevant? (first stage)
        \item Why is it excludable?
        \item What would violate exclusion?
      \end{itemize}
    \end{column}
  \end{columns}

  \vspace{8pt}

  For every method: state the assumption, argue for it, and acknowledge what could go wrong.
\end{frame}

\begin{frame}{Common Empirical Strategy Mistakes}
  \begin{itemize}
    \item \textbf{No equation.} ``I run a regression of Y on X with controls'' isn't enough. Write it out:
    \[Y_{it} = \alpha + \beta D_{it} + X_{it}'\gamma + \mu_i + \lambda_t + \varepsilon_{it}\]

    \vspace{4pt}

    \item \textbf{Skipping the ``why this works'' argument.} Don't just describe what you do. Explain \emph{why} it gives you a credible estimate.

    \vspace{4pt}

    \item \textbf{Mixing in results.} This section is about \emph{design}, not findings. Save the coefficients for the results section.

    \vspace{4pt}

    \item \textbf{No discussion of threats.} Every strategy has weaknesses. Name them. The reader will think of them anyway --- better if you get there first.
  \end{itemize}
\end{frame}

% ══════════════════════════════════════════════════════════
\section{Good vs.\ Bad Examples}
% ══════════════════════════════════════════════════════════

\begin{frame}{Intro: Weak vs.\ Strong}
  \small
  \begin{columns}[T]
    \begin{column}{0.48\textwidth}
      \textbf{\textcolor{red!60}{Weak:}}

      \vspace{2pt}

      ``Minimum wages have been debated for a long time. Many economists have studied this topic. Some find positive effects and some find negative effects. In this paper, I look at minimum wages and employment.''

      \vspace{4pt}

      {\scriptsize \textcolor{darkgray}{Problems: no hook, no stakes, no specific question, vague contribution.}}
    \end{column}
    \begin{column}{0.48\textwidth}
      \textbf{\textcolor{usfgreen}{Strong:}}

      \vspace{2pt}

      ``Twenty-three states raised their minimum wage in 2024, affecting over 8 million workers. Whether these increases reduce employment remains contested. I exploit the staggered timing of state-level increases to estimate the effect on teen employment using a stacked DiD design.''

      \vspace{4pt}

      {\scriptsize \textcolor{darkgray}{Hook (fact), stakes (scale), question (implicit), method (clear).}}
    \end{column}
  \end{columns}
\end{frame}

\begin{frame}{Empirical Strategy: Weak vs.\ Strong}
  \footnotesize
  \begin{columns}[T]
    \begin{column}{0.48\textwidth}
      \textbf{\textcolor{red!60}{Weak:}}

      \vspace{2pt}

      ``I use difference-in-differences. The treatment group is states that raised the minimum wage. The control group is states that didn't. I control for GDP and population.''

      \vspace{4pt}

      {\scriptsize \textcolor{darkgray}{No equation, no identifying assumption, no parallel trends discussion.}}
    \end{column}
    \begin{column}{0.48\textwidth}
      \textbf{\textcolor{usfgreen}{Strong:}}

      \vspace{2pt}

      ``I estimate $Y_{st} = \beta D_{st} + X_{st}'\gamma + \mu_s + \lambda_t + \varepsilon_{st}$, where $D_{st}$ indicates state $s$ raised its minimum wage by period $t$. Identification requires parallel employment trends absent the policy. I provide evidence in Section~5.''

      \vspace{4pt}

      {\scriptsize \textcolor{darkgray}{Equation, named assumption, roadmap to evidence.}}
    \end{column}
  \end{columns}
\end{frame}

% ══════════════════════════════════════════════════════════
\section{GitHub Pages}
% ══════════════════════════════════════════════════════════

\begin{frame}{Not Writing a Traditional Paper?}
  Some of you are presenting your capstone as a blog post, a series of posts, or a website. That's great --- the same principles apply to the writing itself.

  \vspace{6pt}

  You still need:
  \begin{itemize}
    \item A clear question and why it matters (your ``intro'')
    \item Context on what others have found (your ``background'')
    \item A transparent description of your method (your ``empirical strategy'')
    \item Results with honest interpretation
  \end{itemize}

  \vspace{6pt}

  The difference is the format, not the substance. Let's talk about how to set that up.
\end{frame}

\begin{frame}{What Is GitHub Pages?}
  Free static website hosting, straight from a GitHub repo.

  \vspace{6pt}

  \begin{itemize}
    \item Your repo already has code and data. GitHub Pages lets you add a website on top of it.
    \item URL: \texttt{yourusername.github.io/your-repo/}
    \item Supports HTML, Markdown, and --- most usefully --- Quarto.
    \item No server to manage. Push your files, the site updates.
  \end{itemize}

  \vspace{8pt}

  \begin{block}{Two Options}
    \begin{enumerate}
      \item \textbf{Simple:} Write Markdown files, enable Pages, done.
      \item \textbf{Polished:} Use Quarto to build a blog or multi-page site.
    \end{enumerate}
  \end{block}
\end{frame}

\begin{frame}[fragile]{Option 1: Simple Markdown Site}
  \small
  If you just want a single-page write-up:

  \vspace{2pt}

  \begin{enumerate}\setlength\itemsep{0pt}
    \item Create \texttt{docs/index.md} in your repo
    \item Write your content in Markdown
    \item \textbf{Settings} $\to$ \textbf{Pages} $\to$ source: \texttt{main} branch, \texttt{/docs} folder
    \item Push. Your site is live.
  \end{enumerate}

  \vspace{4pt}

  \begin{lstlisting}[style=terminal,basicstyle=\ttfamily\scriptsize]
your-repo/
  docs/
    index.md      # your write-up
    figures/       # images for the site
  code/
    analysis.R
  \end{lstlisting}

  \vspace{2pt}

  GitHub renders the Markdown as HTML automatically. No build step needed.
\end{frame}

\begin{frame}[fragile]{Option 2: Quarto Blog}
  \small
  For multiple posts, a polished layout, or interactive figures:

  \vspace{2pt}

  \begin{lstlisting}[style=terminal,basicstyle=\ttfamily\scriptsize]
# Create a blog project
quarto create project blog my-capstone-blog

# Structure:
my-capstone-blog/
  _quarto.yml           # site config
  index.qmd             # landing page
  posts/
    01-question/
      index.qmd         # "Why I'm studying this"
    02-data/
      index.qmd         # "The data"
    03-results/
      index.qmd         # "What I found"
  \end{lstlisting}

  \vspace{2pt}

  Run \texttt{quarto render} to build, then push the output to GitHub Pages.
\end{frame}

\begin{frame}[fragile]{Setting Up GitHub Pages}
  \small
  \begin{enumerate}\setlength\itemsep{1pt}
    \item Go to your repo on GitHub
    \item \textbf{Settings} $\to$ \textbf{Pages} (left sidebar)
    \item Source: \textbf{Deploy from a branch}, branch: \texttt{main}, folder: \texttt{/docs}
    \item Click \textbf{Save}
  \end{enumerate}

  \vspace{4pt}

  \textbf{For Quarto:} set the output directory in \texttt{\_quarto.yml}:

  \begin{lstlisting}[style=terminal,basicstyle=\ttfamily\scriptsize]
project:
  type: website
  output-dir: docs
  \end{lstlisting}

  \vspace{2pt}

  Then \texttt{quarto render \&\& git add docs/ \&\& git push}. Site updates in about a minute.
\end{frame}

% ══════════════════════════════════════════════════════════
\section{Work Sprint}
% ══════════════════════════════════════════════════════════

\begin{frame}{Work Sprint (6:05--7:30)}
  Pick your track:

  \vspace{8pt}

  \begin{columns}[T]
    \begin{column}{0.48\textwidth}
      \textbf{\textcolor{usfgreen}{Paper Track}}
      \begin{enumerate}
        \item Write your intro (at least 2 paragraphs)
        \item Write your empirical strategy section (at least 1 page)
        \item If time: start the background
      \end{enumerate}
    \end{column}
    \begin{column}{0.48\textwidth}
      \textbf{\textcolor{usfgreen}{Website/Blog Track}}
      \begin{enumerate}
        \item Set up GitHub Pages on your repo
        \item Create your first page or post
        \item Write the ``question and why it matters'' section
      \end{enumerate}
    \end{column}
  \end{columns}

  \vspace{10pt}

  I'll come around for check-ins. If you're stuck on how to start, try this: write the one-sentence version of your contribution, then expand outward.
\end{frame}

% ══════════════════════════════════════════════════════════
\section{Peer Feedback \& Wrap-up}
% ══════════════════════════════════════════════════════════

\begin{frame}{Peer Feedback (7:30--7:55)}
  Swap your intro (or first blog post) with a neighbor.

  \vspace{8pt}

  \textbf{5 minutes:} Read silently. Then answer:

  \vspace{4pt}

  \begin{enumerate}
    \item Can you state their research question in one sentence?
    \item Is the contribution clear --- do you know what they did and what they found?
    \item Where did you get confused or lose the thread?
  \end{enumerate}

  \vspace{8pt}

  If you can't answer \#1 after reading the intro, that's the most useful feedback you can give.
\end{frame}

\begin{frame}{Next Week \& Deadlines}
  \textbf{Next week (4/16):} Figures, tables, slide design, and presentation prep.

  \vspace{10pt}

  \textbf{Upcoming:}
  \begin{itemize}
    \item Weekly progress report due \textbf{Wednesday}
    \item \textbf{\textcolor{usfgreen}{Full Draft due April 24}} (two weeks)
  \end{itemize}

  \vspace{10pt}

  \begin{block}{Between Now and the Full Draft}
    You need a complete paper or website: intro, background, empirical strategy, results, conclusion. Leave today with a drafted intro and empirical strategy, and you're in good shape.
  \end{block}
\end{frame}

\end{document}
